% Folder 147, batch 01 (archivists' pages 1-20).
% Pass: Opus 5.5 (claude-opus-5-5), 2026-10-09 — first pass, unchecked against the pages by a human.
%
% Page 1 is a pencil cover sheet (« Pages 1-90 + table des matières »),
% not in Grothendieck's hand as far as can be judged; it is skipped.
% Page 2 is his table of contents for the whole run (his pp. 1-90).
% Pages 3-20 are his pp. 1-9 of « Structure à l'infini des M_{g,nu} ».
\documentclass[11pt,a4paper]{article}
\input{../preamble/grothendieck}

\folder{147}
\batch{1}
\pages{1}{20}
\foldertitle{Structure à l'infini des Mg,ν [pages 1 à 90, dont table des matières] : notes manuscrites (s.d.).}
\dating{s.d. — le groupe « Autour de La "Longue Marche" à travers la théorie de Galois » (141 à 149) est daté [à partir de 1978]-1983}
\watermark{Édition de démonstration}

\begin{document}

\section{Table des matières}

\page{2}
\note{table des matières de l'auteur pour l'ensemble de la liasse (ses pages 1 à 90) ; les numéros à droite sont ses numéros de page.}

\emph{Structure à l'\uncertain{$\infty$} des $M_{g,\nu}$}

\begin{enumerate}
  \item Courbes standard — 1
  \item Graphe associé à une courbe standard. Maquettes — 3
  \item Courbes stables, et MD-graphes — 6
  \item La théorie de Mumford-Deligne — 10
  \item Spécialisation des MD-graphes — 15
  \item Morphismes de généralisation de graphes et de maquettes — 16'
  \item Étude des $\widehat{M}_{G}$ de dim $\leq 2$ : a) détermination des graphes correspondants — 25
  \item Structure de spécialisation des strates en dim. modulaire ambiante $\leq 4$ — 29'
  \item Structure groupoïdale des multiplicités modulaires de Teichmüller variables (la MDT-structure) : cas transcendant, morphismes de généralisation — 37
  \item Structure MDT analytique : calcul des morphismes de généralisation entre MD-graphes — 45
  \item Digression : déploiement d'un diviseur à croisements normaux. Structure à l'infini des groupoïdes fondamentaux — 52
  \item Digression (suite) : types canoniques associés à une stratification par croisements normaux, et leurs dévissages en types \uncertain{élémentaires} — 74
  \item Digression sur les stratifications \add{« locales »} générales — 86
\end{enumerate}

\section{Structure à l'infini des $M_{g,\nu}$}

\page{3}
\note{p. 1 de l'auteur.}

\subsection{1. Courbes standard}

Soit $k$ un corps alg. clos. Une « courbe standard » sur $k$ est \struck{\ill{}} schéma $X$ sur $k$ satisfaisant les conditions suivantes :

a) $X$ quasi-projectif, \add{toute composante irréd. est} de dim $1$

b) Tout point de $X$ est soit lisse, soit un « point quadratique (ordinaire) » — i.e. isom. (loc. ét) à la courbe $\operatorname{Spec}(k[X,Y]/XY)$ au point $0$.

Il est connu qu'on peut trouver une unique « complétée » $\widehat{X}$ de $X$, telle que \add{$\widehat{X}$ soit un schéma propre, que} $X$ s'identifie à un ouvert dense de $\widehat{X}$, et que $\widehat{X}$ soit lisse en les pts de $\widehat{X} \setminus X = I$. Alors $\widehat{X}$ est une courbe projective, $I$ est \struck{contenu} une partie finie de $\widehat{X}(k)$ contenue dans l'ens. ouvert des pts de lissité de $\widehat{X}$. L'ens. $A$ des pts \emph{singuliers} de $X$ s'identifie à l'ens. \emph{analogue} pour $\widehat{X}$
\begin{equation*}
(1)\quad
\begin{cases}
I = \widehat{X} \setminus X \subset \widehat{X}^{\mathrm{lisse}}(k) & \text{« pts à l'infini » de } X\\
A = X \setminus X^{\mathrm{lisse}} = \widehat{X} \setminus \widehat{X}^{\mathrm{lisse}} \subset X(k) & \text{(pts singuliers de } X\text{)}
\end{cases}
\end{equation*}
La donnée de $X$ équivaut à celle du

\page{4}
$(\widehat{X}, I)$, où $\widehat{X}$ est un schéma projectif, dont les composantes irréductibles sont de dim $1$, et dont l'ensemble singulier est formé de pts quadratiques ordinaires — et $I$ est un sous-schéma fini étale de $\widehat{X}^{\mathrm{lisse}}$, ou ce qui revient au même, une partie finie de $\widehat{X}^{\mathrm{lisse}}(k)$.

Soit $\widehat{\widetilde{X}} = \widetilde{\widehat{X}}$ la normalisée de $\widehat{X}$, \struck{alors}
\[
\widetilde{I} = \widehat{\widetilde{X}} \setminus \widetilde{X} = \text{Image inverse de } I \text{ par } \widehat{\widetilde{X}} \to \widehat{X}
\]
l'\uncertain{ens.} des \struck{ses} « points à l'infini » de $\widetilde{X}$, on a
\[
(2)\qquad \widetilde{I} \xrightarrow{\ \sim\ } I
\]
et nous identifierons généralement $\widetilde{I}$ à $I$, et plus généralement \struck{\ill{}} les pts de $\widehat{X}^{\mathrm{lisse}}$ aux pts correspondants de $\widehat{\widetilde{X}}$. Soit $\widetilde{A}$ l'image inverse de $A$ dans $\widetilde{X}$ (ou dans $\widehat{\widetilde{X}}$), d'où
\[
(3)\qquad \widetilde{A} \longrightarrow A
\]
qui est une application de degré $2$, d'où une involution sans pts fixe $\sigma_{\widetilde{A}}$ dans $\widetilde{A}$. Ainsi, à la courbe standard $X$ nous avons associé le système de données

\page{5}
\note{p. 2 de l'auteur.}

suivant :

(4) a) Une courbe projective et lisse $Y$ $(= \widehat{\widetilde{X}})$ \struck{\ill{}}

b) \struck{Deux} parties \add{finies} disjointes $\widetilde{A}$ et $I$ de $Y(k)$ (ou encore, deux sous-schémas finis étales disjoints $\widetilde{A}_{k}$ et $I_{k}$ de $Y$)

c) Une involution sans point fixe $\sigma_{\widetilde{A}}$ de $\widetilde{A}$ (qui est donc de cardinal pair).

Inversement, partant des données a) b) c), on construit une courbe standard $X$ en passant au quotient dans $Y \setminus I_{k} \supset \widetilde{A}_{k}$ par l'involution $\sigma$ — i.e. $X$ est universel \add{\uncertain{dans} (Sch))} pour les données $p$ :

\begin{tikzcd}
Y \setminus I_{k} \arrow[r, "p"] & X \\
\widetilde{A}_{k} \arrow[u, "i"] \arrow[ur, dashed] &
\end{tikzcd}

soumise à $(pi)\sigma = pi$.

Ainsi la catégorie des courbes standard sur $k$ \add{(pour les iso)} apparaît comme équivalente à celle des systèmes a) b) c) ci-dessus (pour les iso)…

\emph{NB} On récupère $\widehat{X}$ comme quotient de $Y$ par $\sigma$.

\page{6}
\subsection{Généralisation sur une base quelconque}

Une \emph{courbe standard} sur $S$ (multiplicité schématique, disons) sera définie constructivement en termes d'un système a) b) c) comme ci-dessus, i.e.

(5) a) $Y$, schéma propre et lisse \add{sur $S$,} partout de dim relative $1$ ;

b) $\widetilde{A}$ et $I$, deux sous-schémas de $Y$, finis étales sur $S$, \struck{\ill{}} $\widetilde{A} \cap I = \emptyset$ ;

c) $\sigma_{\widetilde{A}}$ une involution de $\widetilde{A}$, opérant librement.

On construit alors $\widehat{X} =$ \struck{$X$} \add{$Y$}$/\sigma$, contenant $A = \widetilde{A}/\sigma$ et $I$ comme sous-schémas \add{fermés} finis étales sur $S$, et on définit $X = \widehat{X} \setminus I$. On peut montrer que le foncteur
\[
(6)\qquad (Y, \widetilde{A}, I, \sigma_{\widetilde{A}}) \longmapsto X
\]
des systèmes (5) (pour les iso) vers les schémas relatifs \add{(pour les iso)}, est \emph{pleinement fidèle}.

\marginal{faux tel quel !}

\marginal{NB NB ce foncteur \uncertain{n'est} \ill{} \ill{} (6) …}
\note{la note marginale « NB NB … » est écrite en oblique dans la marge inférieure gauche et se poursuit, semble-t-il, dans la marge supérieure gauche de la page suivante ; lecture très incertaine.}

\page{7}
\note{p. 3 de l'auteur.}

\marginal{Le foncteur est pl. fid. au mieux sur une base $S$ réduite — faux p. ex. sur les nbs duaux \add{sur un corps} (si $I \neq \emptyset$). Pour avoir un foncteur pl. fidèle, il faut \uncertain{remplacer} \ill{} \ill{} dans le contexte \uncertain{propre}, i.e. \ill{} $(\widehat{X}, I)$, \ill{} $X = \widehat{X} \setminus I$ …}

\ill{} courbe \add{relative} « standard » sur $S$, tout élément de l'image essentielle. \add{Le} système $($\struck{$X$} $Y, \widetilde{A}, I, \sigma)$ qui lui donne naissance s'appellera le « compactifié normalisé » \add{(relatif)} (\ill{} de \struck{système divisant} le \ill{}, puisque c'est non seulement le schéma relatif $Y$, mais $Y$ avec la structure supplémentaire $\widetilde{A}, I, \sigma_{\widetilde{A}}$…).
\note{dans le système, l'auteur a biffé une première lettre (vraisemblablement $X$) ; la lecture $Y$ est une restitution d'après (6).}

\subsection{2. Graphe associé à une courbe standard}

Revenons au cas d'un corps de base $k$ alg. clos, pour commencer. Soit \struck{$X$} une courbe standard, d'où $Y$, \struck{\ill{}} $I$, $\widetilde{A}$, $\sigma_{\widetilde{A}}$.

Posons
\[
(7)\qquad S = \pi_{0}(Y) \simeq \text{ens. des comp. irréd. de } X
\]
On a alors le diagramme d'applications canoniques entre ensembles finis

\begin{tikzcd}
 & I \arrow[d, "p"] \\
\widetilde{A} \arrow[r, "\sigma"] \arrow[d, "q"'] & S \\
A = \widetilde{A}/\sigma &
\end{tikzcd}

(8)
\note{l'application $\widetilde{A} \to S$ est notée $\sigma$, comme l'involution ; ainsi sur la page.}

\page{8}
où $q$ est de degré $2$ et définit l'involution $\sigma_{\widetilde{A}}$. Les applications $\sigma$ et $p$ sont induites par les inclusions $\widetilde{A} \hookrightarrow Y$, $I \hookrightarrow Y$ en passant aux $\pi_{0}$.

Le système $(\widetilde{A} \xrightarrow{\sigma} S, \sigma_{\widetilde{A}})$ ou

\begin{tikzcd}
\widetilde{A} \arrow[r, "\sigma"] \arrow[d, "\deg 2"'] & S \\
A &
\end{tikzcd}

peut être considéré comme définissant un \emph{graphe}, dont $S$ est l'ens. des sommets, et $\widetilde{A}$ l'ens. des \emph{arcs} (ou arêtes \uncertain{munies d'une} orientation) l'application $\sigma$ étant l'application « origine d'un arc ». Ce graphe ne dépend que de $\widehat{X}$, pas de $X$ i.e. du choix de $I \subset \widehat{X}(k)$. C'est pour tenir compte de ce choix que l'on considère, en plus de la structure de graphe, la donnée supplémentaire
\[
(9)\qquad I \longrightarrow S
\]

\page{9}
\note{p. 4 de l'auteur.}

Le graphe indique comment les composantes irréductibles de $X$ (figurées par les sommets) se recoupent deux à deux — les pts d'intersection \add{« doubles », i.e.} i.e. les pts singuliers (quadratiques) de $X$, correspondant aux arêtes. Si une composante irréductible $X_{\alpha}$ de $X$ correspond au sommet $\alpha$ du graphe, alors les boucles fermées en $\alpha$ correspondent bijectivement aux pts doubles de $X_{\alpha}$ — donc dire que les $X_{\alpha}$ sont lisses signifie que le graphe $G_{X}$ n'a pas de boucle fermée, i.e. que pour tout arc, l'origine est $\neq$ de l'extrémité.

Il est clair que tout graphe fini peut être obtenu \add{(à iso près)} \struck{à l'aide} par un $\widehat{X}$ convenable — et \ill{} \ill{} des compo-

\page{10}
santes $X_{\alpha}$ de genre \add{$g_{\alpha}$} donné (i.e. des $\widetilde{X}_{\alpha}$ de genre $g_{\alpha}$ …). De plus, quand on se donne \struck{des} $I \to S$ ($I$ un fini), cela peut être réalisé par un $I \subset \widehat{X}^{\mathrm{lisse}}$, i.e. par une courbe standard \uncertain{$S$}.
\note{on attendrait $X$.}

La \emph{maquette} d'une courbe standard $X$ consiste, par définition, en les données suivantes

(10) a) Le graphe $G_{X} = (S, \widetilde{A}, \sigma_{\widetilde{A}})$, où $S = \pi_{0}(\widetilde{X}) = \pi_{0}(\widehat{\widetilde{X}})$, $\widetilde{A} =$ \uncertain{$\operatorname{Sing}(X)$} \struck{\ill{}}
\note{les égalités sont écrites en colonne sous les termes du triplet ; la lecture « Sing$(X)$ » sous $\widetilde{A}$ est douteuse.}

b) L'\struck{application} \add{ensemble} $I = \widehat{X} \setminus X$, et l'application
\[
I \longrightarrow S
\]

c) L'application
\[
S \xrightarrow{\ g\ } \mathbb{N}, \qquad \alpha \longmapsto \text{genre de } Y_{\alpha} = \widehat{\widetilde{X}}_{\alpha}.
\]

Une structure \add{formée} d'un graphe fini $G$ \add{$= (S, \widetilde{A}, \sigma_{\widetilde{A}})$}, d'un ensemble fini $I$ \struck{$A$} au-dessus de l'ens. des sommets de $G$, et d'une application « genre » : $S \xrightarrow{g} \mathbb{N}$, sera

\page{11}
\note{p. 5 de l'auteur.}

appelée ici une « maquette ».

\add{Considérons la maquette d'une courbe standard $X$.}

\emph{Proposition} a) Soient $\alpha, \beta \in S$, alors $\alpha, \beta$ appartiennent à la même composante connexe du graphe $G$, ssi $X_{\alpha}$ et $X_{\beta}$ appartiennent à la même comp. conn. de $X$. Donc \struck{\ill{}} on a une bij. can.
\[
(11)\qquad \pi_{0}(G_{X}) \simeq \pi_{0}(X),
\]
en particulier $X$ est connexe ssi $G_{X}$ est connexe.

b) Supposons $X$ connexe i.e. $\widehat{X}$ connexe, i.e. $H^{0}(\widehat{X}, \mathcal{O}_{\widehat{X}}) \simeq k$ i.e. $\dim H^{0}(\widehat{X}, \mathcal{O}_{\widehat{X}}) = 1$.

Posons \struck{$g(X) = \dim H^{1}(X, \mathcal{O}_{X})$, $\chi(X) = \chi$}
\begin{equation*}
(12)\quad
\begin{cases}
g(\widehat{X}) = \dim H^{1}(\widehat{X}, \mathcal{O}_{\widehat{X}})\\
\chi_{\mathrm{coh}}(\widehat{X}) = \chi(X, \mathcal{O}_{X}) = 1 - g(X)
\end{cases}
\end{equation*}
\note{dans le second membre de la deuxième ligne, un premier argument de $\chi$ est surchargé et biffé ; $X$ est lu à la place de $\widehat{X}$, tel quel.}

on a alors
\[
(13)\qquad \chi_{\mathrm{coh}}(\widehat{X}) = \chi(\widehat{\widetilde{X}}) - \underbrace{\operatorname{card} A}_{\mu} = \sum_{\alpha \in S} (1 - g_{\alpha}) - \mu
\]
\note{un premier second membre, $\sum_{\alpha \in S} \chi_{\mathrm{coh}}(\ldots)$, est biffé et surchargé ; la forme retenue est écrite au-dessus.}

i.e.
\begin{equation*}
(14)\quad
\begin{aligned}
g(\widehat{X}) &= \sum g_{\alpha} + (-\operatorname{card} S + \underbrace{\operatorname{card} A}_{\mu} + 1)\\
&= \sum g_{\alpha} + h_{1}
\end{aligned}
\end{equation*}

\page{12}
où \struck{\ill{}}
\[
(15)\qquad h_{1} = \operatorname{rg}_{\mathbb{Z}} H_{1}(G_{X}, \mathbb{Z}) .
\]

\emph{Démonstration} Standard (!) — la formule (13) résulte d'une suite exacte \add{standard} de cohomologie reliant $\mathcal{O}_{\widehat{X}}$, $\mathcal{O}_{\widehat{\widetilde{X}}}$ et les singularités de $\widehat{X}$, d'ailleurs la formule d'EP pour le graphe $G_{X}$ nous donne
\[
\chi(G_{X}) = \underbrace{h_{0}}_{1} - h_{1} = \operatorname{card} S - \operatorname{card} A
\]
i.e.
\[
1 - \operatorname{card} S + \operatorname{card} A = h_{1}, \qquad \text{O.K.}
\]

\emph{Corollaire} Le nb $g(\widehat{X})$ ne dépend que de la maquette de $X$, par la formule (14)
\[
g(\widehat{X}) = \underbrace{\sum_{\alpha \in S} g(\alpha) + h_{1}}_{\overset{\mathrm{d\acute ef}}{=}\ \text{genre de } G}
\]
\ill{} appelle \struck{le genre} \add{genre de $G$} de la maquette.
\note{l'accolade sous la formule porte « $\overset{\mathrm{déf}}{=}$ genre de $G$ » ; la ligne qui suit est corrigée et mal lisible.}

On appelle \emph{type} de la maquette le couple $(g, \nu)$, où $g$ est son genre et $\nu = \operatorname{card} I$ est le \emph{poids marqué}. On appelle \emph{poids marqué} d'un sommet l'entier

\page{13}
\note{p. 6 de l'auteur.}

\marginal{\ill{} \ill{} \ill{}}
\ill{} $\operatorname{card} I_{\alpha}$ (où $I_{\alpha}$ est la fibre de $I$ en $\alpha \in S$)
\[
(16)\qquad \nu_{\alpha} = \operatorname{card} I_{\alpha} ,
\]
et on appelle \emph{poids total} d'un sommet l'entier
\[
(17)\qquad \widehat{\nu}_{\alpha} = \nu_{\alpha} + \omega_{\alpha}
\]
$\omega_{\alpha} \overset{\mathrm{d\acute ef}}{=}$ ordre du sommet $\alpha$, i.e. $\operatorname{card} \widetilde{A}_{\alpha}$, où $\widetilde{A}_{\alpha}$ est la fibre de $\widetilde{A}$ au-dessus de $\alpha$, i.e. le nb d'arcs sortant de $\alpha$. C'est donc le nb de « pts marqués » (provenant de pts doubles ou de pts à l'infini de $X$) sur la composante $Y_{\alpha}$ de $Y = \widehat{\widetilde{X}}$.
\note{une courte note oblique dans la marge gauche, en haut de la page, n'est pas lue.}

\subsection{3. Courbes « stables » et MD-graphes}

Une courbe standard \add{(sur $k$ alg. clos)} est « stable », si elle satisfait à l'une des conditions équivalentes suivantes

a) $\operatorname{Aut} X$ est fini

b) Pour tout $\alpha$, $(Y_{\alpha}, I_{\alpha} \cup \widetilde{A}_{\alpha})$ est anabélienne i.e. $2g_{\alpha} + \widehat{\nu}_{\alpha} \geq 3$ i.e. $2g_{\alpha} - 2 + \widehat{\nu}_{\alpha} \geq 1$, i.e.

\page{14}
1°) si $g_{\alpha} = 1$, on a $\widehat{\nu}_{\alpha} \geq 1$

2°) si $g_{\alpha} = 0$ on a $\widehat{\nu}_{\alpha} \geq 3$

c) Tout champ de vecteurs sur $Y$ nul sur \struck{\ill{}} $I \cup \widetilde{A}$ est nul.

d) $\underline{\operatorname{Aut}}_{(Y, \widetilde{A}_{k}, I)}$ est un schéma en groupes fini étale \add{sur $k$}.

On voit que cette condition (sous forme b)) ne dépend que de la \emph{maquette} de la courbe. On dit que $X$ est une \emph{MD-courbe} (MD, initiales de « Mumford-Deligne » — ou « modulaire ») si elle est stable, et \add{connexe} (i.e. connexe non vide). Les maquettes de telles courbes sont les maquettes \add{connexes et} stables (i.e. dont les sommets de genre $1$ sont de poids total $\geq 1$, et les sommets de genre $0$ sont de poids total $\geq 3$), on les appellera les MD-graphes.

\emph{NB} Une maquette est un MD-graphe ssi

a) elle est connexe (i.e. le graphe $G$ est conn. $\neq \emptyset$)

b) elle n'est pas réduite à un seul sommet de genre $1$, de poids total $0$ (pas d'arête, pas de points marqués)

c) les sommets de genre $0$ sont de poids total $\geq 3$.

\page{15}
\note{p. 7 de l'auteur.}

\emph{Proposition} Si $(G = (S, \widetilde{A}, \sigma), I, g : S \to \mathbb{N})$ est un MD-graphe, son type $(g, \nu)$ est anabélien, i.e. $2g + \nu \geq 3$.

Si on avait $g = 1$, $\nu = 0$, alors la relation
\[
g = 1 = \sum g_{\alpha} + h_{1}
\]
montre que ou bien tous les $g_{\alpha}$ sont nuls et $h_{1} =$ \struck{$1$} $1$ \add{(i.e. \uncertain{$G$ est un cercle})}, ou bien tous les $g_{\alpha}$ sauf un \add{$g_{\alpha_{0}}$} sont nuls, et $g_{\alpha_{0}} = 1$, $h_{1} = 0$. \add{Dans le premier cas,} (La relation $s_{0} - s_{1} = h_{0} - h_{1} = 1 - 1 = 0$ montre qu'il y a autant de sommets que d'arêtes, dans le deuxième, $s_{0} - s_{1} = 1 - 0 = 1$ il y a un sommet de plus que d'arêtes.)

Mais la relation $\nu = 0$ implique que les sommets de genre $0$ sont d'ordre $\geq 3$. \add{($G$ n'étant pas réduit à un sommet, puisqu'il est stable d'après \ill{})} \struck{Mais un} \add{arbre} qui n'est réduit à un seul sommet \struck{graphe top. connexe ou} a au moins deux bouts, dont l'un au moins sera de genre $0$, ce qui élimine le cas 2°. Donc tous les sommets sont de genre $0$ \add{(donc aucun sommet n'est un bout)} et $h_{1} = 1$. Mais \struck{comme} un graphe \add{topologique} $0$-connexe avec $h_{1} = 1$ \add{pt.} \uncertain{bout} topologique est homéomorphe \uncertain{à un} cercle, donc s'il est réalisé par un graphe combinatoire, tous les sommets de celui-ci sont d'ordre $2$, ce qui contredit qu'ils soient d'ordre trois.
\note{page très corrigée : plusieurs ajouts interlinéaires se chevauchent ; l'ordre de lecture retenu est conjectural.}

\page{16}
Éliminons de même le cas $g = 0$, $\nu \leq 2$. Ici
\[
g = 0 = \sum g_{\alpha} + h_{1}
\]
implique que les $g_{\alpha}$ et $h_{1}$ sont nuls : tous les sommets sont de genre $0$, le graphe est un arbre. Il ne peut être réduit à un seul sommet, donc il a au moins deux bouts, donc \struck{\ill{}} les \struck{poids} \uncertain{poids} marqués doivent être $\geq 2$, donc \ill{}
\[
\nu \geq 2 + 2 = 4, \quad \text{absurde.}
\]

Soit $G$ une \struck{MD-graphe} \add{maquette}. On dit qu'une \struck{MD} \add{standard} courbe \struck{standard} sur un corps alg. clos est \emph{de type $G$}, si sa maquette est iso à $G$, on dit qu'elle est \add{$G$-}\emph{épinglée} si on s'est donné un iso entre sa maquette et $G$ (c'est donc une structure supplémentaire sur $X$, appelée $G$-épinglage).

Soit $(\widehat{X}, \underline{I})$ une courbe standard sur une base $S$ quelconque, on dit qu'elle est de type $G$ si ses fibres géom. sont de type $G$. Alors les \struck{\ill{}} \add{maquettes} des

\page{17}
\note{p. 8 de l'auteur.}

fibres géom. de $(\widehat{X}, \underline{I})$ forment les fibres d'un schéma en maquettes \add{sur $S$ (ou en MD-graphes)}
\[
(\underline{S}, \underline{\widetilde{A}}, \sigma_{\underline{\widetilde{A}}}, \underline{I}, \underline{\widetilde{A}} \to \underline{S}, \underline{I} \to \underline{S}, \underline{S} \xrightarrow{g} \mathbb{N}_{S})
\]
(système de rev. finis étales \add{de $S$} et de morphismes entre ceux-ci), loc. isomorphe à la maquette $G$ donnée. \struck{\ill{}} On appelle \add{$G$-}épinglage de $(\widehat{X}, \underline{I})$ tout isomorphisme entre $G_{S}$ et $\underline{G}(\widehat{X}, \underline{I})$. Si \struck{$\Gamma$}
\[
(18)\qquad \Gamma = \operatorname{Aut} G
\]
(groupe fini), les $G$-épinglages de $(\widehat{X}, \underline{I})$ s'identifient aux sections d'un certain $\Gamma_{S}$-torseur, \struck{qui est le} \add{appelé} \emph{torseur des $G$-épinglages} de $(\widehat{X}, \underline{I})$.

Considérons, sur une base $\mathcal{S}$ fixe, la catégorie \add{(pour les iso)} des courbes standard $G$-épinglées. Pour tout $\alpha \in S$ (ens. des sommets de la maquette) soit $\widehat{I}_{\alpha} = I_{\alpha} \amalg \widetilde{A}_{\alpha}$ \uncertain{l'élément} (de cardinal $\widehat{\nu}_{\alpha}$ = poids total de $\alpha$) somme

\page{18}
des fibres $I_{\alpha}$ et $\widetilde{A}_{\alpha}$ de $I$ et $A$ au-dessus de $\alpha$. La donnée d'une courbe standard $G$-épinglée sur $\mathcal{S}$ équivaut à la donnée, pour tout $\alpha \in S$, d'une courbe relative (propre lisse \add{relt. connexe}) $Y_{\alpha}$ sur $\mathcal{S}$, de genre relatif $g_{\alpha}$, \emph{et} d'un \add{$\mathcal{S}$-}monomorphisme $(\widehat{I}_{\alpha})_{\mathcal{S}} \hookrightarrow Y_{\alpha}$.

Désignons, \struck{\ill{}} pour \struck{deux} \add{un} entier naturel $g$ et un ens. fini $I$ (de cardinal $\nu$) par $\mathcal{M}_{g,I}(\mathcal{S})$ le \struck{catégorie} \add{groupoïde} (19) des courbes relatives $Y$ (propres, lisses, connexes) sur $\mathcal{S}$, $I$-\emph{ponctuées} i.e. munies d'un mono $I_{\mathcal{S}} \hookrightarrow Y$. On \uncertain{désignera} de même,
\[
(20)\qquad \mathcal{M}_{G}(\mathcal{S})
\]
désigne le groupoïde des courbes standard $G$-épinglées sur $\mathcal{S}$, on a une équivalence
\[
(21)\qquad \mathcal{M}_{G}(\mathcal{S}) \approx \prod_{\alpha \in S} \mathcal{M}_{g_{\alpha}, \widehat{I}_{\alpha}}(\mathcal{S}) .
\]
\note{$\mathcal{M}_{g,I}(\mathcal{S})$ est entouré et relié par une flèche au numéro (19) placé devant « des courbes ». Dans la deuxième phrase, une formule biffée (une première notation $\mathcal{M}_{g,\ldots}$) n'est pas lue.}

\page{19}
\note{p. 9 de l'auteur.}

Cette équivalence est compatible avec les chgts de base, et avec le transport de structure pour $G$ variable.

On voit donc que, pour que, pour $\mathcal{S}$ variable, le champ en groupoïdes
\[
(22)\qquad \mathcal{S} \longmapsto \mathcal{M}_{G}(\mathcal{S})
\]
soit $2$-représentable par une étendue schématique, il f. et s. qu'il en soit de même des
\[
\mathcal{S} \longmapsto \mathcal{M}_{g_{\alpha}, \widehat{I}_{\alpha}}(\mathcal{S}) ,
\]
ce qui signifie aussi que les types correspondants $(g_{\alpha}, \widehat{\nu}_{\alpha})$ sont anabéliens, i.e. que la maquette $G$ \add{donnée} est stable. Si on suppose $G$ connexe, cela signifie donc justement que $G$ est un MD-graphe. On désignera alors par $\mathcal{M}_{G}$ la multiplicité modulaire

\page{20}
correspondante, qui n'est qu'à un iso can. près

\marginal{\ill{}, une équivalence}
\[
(23)\qquad \mathcal{M}_{G} \simeq \prod_{\alpha} \mathcal{M}_{g_{\alpha}, \widehat{I}_{\alpha}} .
\]

Notons que $\Gamma = \operatorname{Aut} G$ opère sur $\mathcal{M}_{G}$ de façon évidente, et aussi sur le système $(\widehat{I}_{\alpha}, g_{\alpha})_{\alpha \in S}$ donc sur le deuxième membre de (23), qui \struck{est} \add{est un} iso. compatible aux op. de $\Gamma$.

On peut considérer la multiplicité mixte
\[
(24)\qquad (\mathcal{M}_{G}, \Gamma) = \mathcal{M}_{[G]} ,
\]
qui a priori ne dépend que du type d'isomorphie $[G]$ de $G$, et qui s'interprète, en termes de ce dernier, comme la multiplicité schématique modulaire pour le groupoïde fibré (sur la $2$-cat. des multiplicités schématiques) des courbes standard relatives de type $[G]$.

\note{le texte se poursuit au-delà de ce lot (p. 10 de l'auteur, « 4. La théorie de Mumford-Deligne » selon la table des matières).}

\end{document}
